%% CerT-MCMC V2 Theory Note — Quantile-Core Certificate
%% Draft for internal use; to be integrated into V2 manuscript
%% Jun Hu, May 2026

\documentclass[11pt]{article}
\usepackage{amsmath,amssymb,amsthm}
\usepackage[margin=2.5cm]{geometry}
\usepackage{enumitem}

\newtheorem{theorem}{Theorem}
\newtheorem{proposition}[theorem]{Proposition}
\newtheorem{corollary}[theorem]{Corollary}
\newtheorem{lemma}[theorem]{Lemma}
\newtheorem{remark}[theorem]{Remark}
\newtheorem{definition}[theorem]{Definition}
\newtheorem{assumption}{Assumption}

\newcommand{\R}{\mathbb{R}}
\newcommand{\E}{\mathbb{E}}
\newcommand{\Prob}{\mathbb{P}}
\newcommand{\osc}{\mathrm{osc}}
\newcommand{\KL}{\mathrm{KL}}

\title{CerT-MCMC V2: Quantile-Core Spectral Gap Certificate\\[4pt]
\large Theory Note (Draft v0.1)}
\author{Jun Hu}
\date{May 2026}

\begin{document}
\maketitle

%% ═══════════════════════════════════════════════════════
\section{Setup and Notation}
%% ═══════════════════════════════════════════════════════

\textbf{Target.} Let $\pi(\theta)$ be the (possibly unnormalised) target density on $\Theta \subseteq \R^D$.

\textbf{Transport map.} Let $T_\phi: \R^D \to \Theta$ be a learned normalising flow with parameters $\phi$, and let $T_\phi^{-1}$ denote its inverse.

\textbf{Reference.} Let $q_0 = \mathcal{N}(0, I_D)$ denote the standard Gaussian in the latent $z$-space.

\textbf{Pullback target.} Define the pullback of $\pi$ to $z$-space:
\[
  \pi_z(z) \;=\; \pi\bigl(T_\phi(z)\bigr) \,\bigl|\det J_{T_\phi}(z)\bigr|.
\]
Both $q_0$ and $\pi_z$ are densities on $\R^D$. The normalising constant of $\pi_z$ equals that of $\pi$.

\textbf{Log importance weight.} Define
\begin{equation}\label{eq:residual}
  r(z) \;=\; \log \frac{\pi_z(z)}{q_0(z)}
  \;=\; \log \pi\bigl(T_\phi(z)\bigr) + \log\bigl|\det J_{T_\phi}(z)\bigr| + \tfrac{1}{2}\|z\|^2 + \tfrac{D}{2}\log(2\pi).
\end{equation}
If $T_\phi$ were a perfect transport map, $r(z)$ would be constant (equal to $\log Z_\pi + \frac{D}{2}\log(2\pi)$, where $Z_\pi = \int \pi(\theta)\,d\theta$). The \emph{oscillation} of $r$ measures the quality of the transport map.

\textbf{Truncated proposal.} Fix $\alpha \in (0,1)$ and let
\[
  K_\alpha = \{z \in \R^D : \|z\| \le R_\alpha\}, \qquad R_\alpha = \sqrt{\chi^2_D(1-\alpha)},
\]
so that $q_0(K_\alpha) = 1-\alpha$. Define $\mu_K = q_0(\cdot \mid K_\alpha)$.

\textbf{Independence Metropolis--Hastings.} The $K_\alpha$-restricted independence MH kernel targets $\pi_z(\cdot \mid K_\alpha)$ with proposal $\mu_K$, accepting $z'$ proposed from state $z$ with probability
\[
  a(z, z') = \min\!\Bigl(1,\; \frac{w(z')}{w(z)}\Bigr), \qquad w(z) = \frac{\pi_z(z)}{q_0(z)} = e^{r(z)}.
\]

\textbf{V1 certificate (covering-based).} The V1 CerT-MCMC certificate bounds
\[
  C_{\mathrm{V1}} = \osc_{K_\alpha}(r) = \sup_{z \in K_\alpha} r(z) - \inf_{z \in K_\alpha} r(z)
\]
using a finite-sample covering argument with correction $O(n^{-1/D})$, yielding the spectral gap bound $\gamma \ge 2/(1+e^{C_{\mathrm{V1}}})$ via Mengersen--Tweedie. The $O(n^{-1/D})$ rate makes V1 certificates vacuous for $D \gtrsim 6$.


%% ═══════════════════════════════════════════════════════
\section{Residual Core}
%% ═══════════════════════════════════════════════════════

The key observation motivating V2 is empirical: for a well-trained flow, the full oscillation $\osc_{K_\alpha}(r)$ is dominated by rare tail outliers, while $r(z)$ is highly concentrated for the bulk of the proposal mass.

\begin{definition}[Residual core]\label{def:core}
For $\rho \in (0, \tfrac{1}{2})$, let $q_\rho$ and $q_{1-\rho}$ denote the $\rho$- and $(1-\rho)$-quantiles of $r(Z)$ under $Z \sim \mu_K$:
\[
  q_\rho = \inf\{t : F_r(t) \ge \rho\}, \qquad q_{1-\rho} = \inf\{t : F_r(t) \ge 1-\rho\},
\]
where $F_r(t) = \mu_K(\{z : r(z) \le t\})$ is the CDF of $r(Z)$. Define the \emph{residual core}
\begin{equation}\label{eq:core}
  G_\rho = \{z \in K_\alpha : q_\rho \le r(z) \le q_{1-\rho}\}.
\end{equation}
\end{definition}

By construction:
\begin{equation}\label{eq:core-mass}
  \mu_K(G_\rho) \ge 1 - 2\rho,
\end{equation}
and the oscillation of $r$ on $G_\rho$ satisfies
\begin{equation}\label{eq:core-osc}
  \osc_{G_\rho}(r) = \sup_{z \in G_\rho} r(z) - \inf_{z \in G_\rho} r(z) \le q_{1-\rho} - q_\rho =: C_\rho.
\end{equation}
The quantity $C_\rho$ depends only on the one-dimensional distribution of $r(Z)$ and is \emph{independent of the ambient dimension~$D$}.


%% ═══════════════════════════════════════════════════════
\section{Main Results}
%% ═══════════════════════════════════════════════════════

% ─── Theorem A ───────────────────────────────────────

\begin{theorem}[Finite-sample quantile certificate]\label{thm:quantile}
Let $z_1, \dots, z_n \overset{\mathrm{iid}}{\sim} \mu_K$ and $r_i = r(z_i)$, $i=1,\dots,n$. Fix confidence level $\zeta \in (0,1)$ and trimming level $\rho \in (0, \tfrac{1}{2})$. Define the DKW bandwidth
\begin{equation}\label{eq:dkw-eps}
  \varepsilon_n = \sqrt{\frac{\log(2/\zeta)}{2n}}.
\end{equation}
Assume $\rho > \varepsilon_n$ (equivalently, $n > \log(2/\zeta)/(2\rho^2)$). Let $\hat{Q}_p$ denote the empirical $p$-quantile of $\{r_1,\dots,r_n\}$. Then with probability at least $1-\zeta$:
\begin{equation}\label{eq:cert-bound}
  C_\rho \;\le\; \hat{C}_\rho \;:=\; \hat{Q}_{1-\rho+\varepsilon_n} - \hat{Q}_{\rho-\varepsilon_n}.
\end{equation}
In particular, $\hat{C}_\rho$ is a valid high-confidence upper bound on the core oscillation $C_\rho$, and its estimation rate is $O(n^{-1/2})$, independent of $D$.
\end{theorem}

\begin{proof}[Proof sketch]
By the Dvoretzky--Kiefer--Wolfowitz (DKW) inequality \cite{DKW1956,Massart1990}, with probability $\ge 1-\zeta$,
\[
  \sup_{t \in \R} |F_n(t) - F_r(t)| \le \varepsilon_n,
\]
where $F_n$ is the empirical CDF of $\{r_i\}$.

On this event, for any $p \in [\varepsilon_n, 1-\varepsilon_n]$:
\begin{align}
  F_r(\hat{Q}_p) &\ge F_n(\hat{Q}_p) - \varepsilon_n \ge p - \varepsilon_n, \label{eq:dkw-lower}\\
  F_r(\hat{Q}_p) &\le F_n(\hat{Q}_p) + \varepsilon_n \le p + \varepsilon_n. \label{eq:dkw-upper}
\end{align}
Therefore, $\hat{Q}_p$ falls between the true $(p-\varepsilon_n)$- and $(p+\varepsilon_n)$-quantiles.

For the upper quantile: $q_{1-\rho} = F_r^{-1}(1-\rho) \le F_r^{-1}(F_r(\hat{Q}_{1-\rho+\varepsilon_n})) \le \hat{Q}_{1-\rho+\varepsilon_n}$, where the first inequality uses \eqref{eq:dkw-lower} with $p = 1-\rho+\varepsilon_n$, giving $F_r(\hat{Q}_{1-\rho+\varepsilon_n}) \ge 1-\rho$.

Similarly, $q_\rho \ge \hat{Q}_{\rho-\varepsilon_n}$, using \eqref{eq:dkw-upper} with $p = \rho-\varepsilon_n$.

Combining: $C_\rho = q_{1-\rho} - q_\rho \le \hat{Q}_{1-\rho+\varepsilon_n} - \hat{Q}_{\rho-\varepsilon_n} = \hat{C}_\rho$.
\end{proof}

\begin{remark}[Dimension-free rate]
The certificate $\hat{C}_\rho$ converges to $C_\rho$ at rate $O(n^{-1/2})$ (via the DKW bandwidth $\varepsilon_n$). This is in stark contrast to V1's covering-based certificate, whose correction term scales as $O(n^{-1/D})$. For $D=20$ and $n=200{,}000$, V1 has $\varepsilon^* \approx 4.2$, while V2 has $\varepsilon_n \approx 0.003$.
\end{remark}

\begin{remark}[Tighter alternative: order-statistic bound]\label{rem:order-stat}
A tighter finite-sample bound replaces the DKW inequality with exact binomial confidence intervals for order statistics. For the $\rho$-quantile, the $k$-th order statistic $r_{(k)}$ satisfies $\Prob(r_{(k)} \le q_\rho) = \Prob(\mathrm{Bin}(n,\rho) \ge k)$. This gives tighter $\hat{C}_\rho$ for moderate $n$ but the same $O(n^{-1/2})$ asymptotic rate. Either version is valid for certification.
\end{remark}

% ─── Theorem B ───────────────────────────────────────

\begin{theorem}[Core-restricted spectral gap]\label{thm:core-gap}
Let $G_\rho$ be the residual core (Definition~\ref{def:core}) and let $P_{G_\rho}$ denote the independence Metropolis--Hastings kernel on $G_\rho$:
\begin{itemize}[nosep]
  \item proposal: $\mu_G = q_0(\cdot \mid G_\rho)$,
  \item target: $\pi_G = \pi_z(\cdot \mid G_\rho)$,
  \item acceptance: $a_G(z,z') = \min\bigl(1,\; w_G(z')/w_G(z)\bigr)$, where $w_G(z) = \pi_G(z)/\mu_G(z)$.
\end{itemize}
Then the spectral gap of $P_{G_\rho}$ in $L^2(\pi_G)$ satisfies
\begin{equation}\label{eq:core-spectral-gap}
  \gamma(P_{G_\rho}) \;\ge\; \frac{2}{1 + e^{C_\rho}},
\end{equation}
where $C_\rho = q_{1-\rho} - q_\rho$ is the quantile oscillation from~\eqref{eq:core-osc}.
\end{theorem}

\begin{proof}[Proof]
By the Mengersen--Tweedie (1996) theorem, the spectral gap of an independence MH kernel with weight function $w$ satisfies
\[
  \gamma \;\ge\; \frac{2}{1 + \exp\!\bigl(\osc(\log w)\bigr)}.
\]
It suffices to show $\osc_{G_\rho}(\log w_G) \le C_\rho$.

Compute:
\[
  \log w_G(z) = \log \frac{\pi_G(z)}{\mu_G(z)}
  = \log \frac{\pi_z(z) / \pi_z(G_\rho)}{q_0(z) / q_0(G_\rho)}
  = \underbrace{\log \frac{\pi_z(z)}{q_0(z)}}_{= r(z)} + \underbrace{\log \frac{q_0(G_\rho)}{\pi_z(G_\rho)}}_{\text{constant in } z}.
\]
Therefore,
\[
  \osc_{G_\rho}(\log w_G) = \osc_{G_\rho}(r) \le q_{1-\rho} - q_\rho = C_\rho,
\]
where the inequality is by~\eqref{eq:core-osc}. Apply Mengersen--Tweedie to $P_{G_\rho}$.
\end{proof}

\begin{remark}[What this theorem certifies]
Theorem~\ref{thm:core-gap} certifies the spectral gap of the independence MH chain \emph{restricted to the residual core $G_\rho$}. It does \textbf{not} directly certify the spectral gap of the full $K_\alpha$-restricted chain. However, the core $G_\rho$ captures at least $(1-2\rho)$ of the proposal mass (equation~\eqref{eq:core-mass}), and typically a comparable fraction of the target mass (Proposition~\ref{prop:target-mass} below). The core certificate therefore characterises the sampler's mixing behaviour on the dominant portion of the posterior.
\end{remark}

% ─── Combining A + B ─────────────────────────────────

\begin{corollary}[Certified core spectral gap from data]\label{cor:certified}
Under the conditions of Theorems~\ref{thm:quantile} and~\ref{thm:core-gap}, with probability at least $1-\zeta$ over the certification sample:
\begin{equation}\label{eq:certified-gap}
  \gamma(P_{G_\rho}) \;\ge\; \frac{2}{1 + e^{\hat{C}_\rho}},
\end{equation}
where $\hat{C}_\rho = \hat{Q}_{1-\rho+\varepsilon_n} - \hat{Q}_{\rho-\varepsilon_n}$ is the empirical quantile certificate from Theorem~\ref{thm:quantile}.
\end{corollary}

\begin{proof}
On the event $\{C_\rho \le \hat{C}_\rho\}$ (probability $\ge 1-\zeta$ by Theorem~\ref{thm:quantile}), apply Theorem~\ref{thm:core-gap} and monotonicity of $x \mapsto 2/(1+e^x)$.
\end{proof}

% ─── Proposition C ───────────────────────────────────

\begin{proposition}[Target mass of the core]\label{prop:target-mass}
Let $\bar{\pi}_z = \pi_z(\cdot \mid K_\alpha)$ denote the target restricted to $K_\alpha$. Then:
\begin{equation}\label{eq:target-mass}
  \bar{\pi}_z(G_\rho) \;\ge\; (1-2\rho)\,\frac{e^{q_\rho}}{e^{q_{1-\rho}}}
  \;=\; (1 - 2\rho)\,e^{-C_\rho}.
\end{equation}
\end{proposition}

\begin{proof}[Proof sketch]
For any measurable $A \subseteq K_\alpha$:
\[
  \bar{\pi}_z(A) = \frac{\int_A \pi_z(z)\,dz}{\int_{K_\alpha} \pi_z(z)\,dz}
  = \frac{\int_A e^{r(z)} q_0(z)\,dz}{\int_{K_\alpha} e^{r(z)} q_0(z)\,dz}.
\]
On $G_\rho$, $r(z) \ge q_\rho$, so
\[
  \int_{G_\rho} e^{r(z)} q_0(z)\,dz \ge e^{q_\rho} \cdot q_0(G_\rho).
\]
On $K_\alpha$, $r(z) \le \sup_{K_\alpha} r \le q_{1-\rho} + \Delta$, but more simply, for an upper bound on the denominator we use
\[
  \int_{K_\alpha} e^{r(z)} q_0(z)\,dz \le e^{\sup_{K_\alpha} r} \cdot q_0(K_\alpha).
\]
This gives a loose bound. A tighter approach: partition $K_\alpha = G_\rho \cup G_\rho^c$.
\begin{align*}
  \bar{\pi}_z(G_\rho) &= \frac{\E_{\mu_K}[e^{r(Z)} \mathbf{1}_{G_\rho}(Z)]}{\E_{\mu_K}[e^{r(Z)}]} \\
  &\ge \frac{e^{q_\rho} \cdot \mu_K(G_\rho)}{e^{q_\rho} \cdot \mu_K(G_\rho) + e^{\sup r} \cdot \mu_K(G_\rho^c)} \\
  &\ge \frac{(1-2\rho) \, e^{q_\rho}}{(1-2\rho) \, e^{q_\rho} + 2\rho \, e^{\sup r}}.
\end{align*}
This still involves $\sup r$, which is the V1 quantity. An alternative that avoids $\sup r$ entirely:
\begin{align}
  \bar{\pi}_z(G_\rho)
  &= 1 - \bar{\pi}_z(G_\rho^c) \nonumber \\
  &= 1 - \frac{\E_{\mu_K}[e^{r(Z)} \mathbf{1}_{G_\rho^c}(Z)]}{\E_{\mu_K}[e^{r(Z)}]}. \label{eq:target-mass-alt}
\end{align}
With a finite-sample estimate of $\E[e^r]$ and $\E[e^r \mathbf{1}_{G_\rho^c}]$, this can be bounded empirically. However, a rigorous finite-sample bound requires care with the exponential; self-normalised importance sampling concentration inequalities may apply.
\end{proof}

\begin{remark}[Practical sufficiency of the proposal-mass statement]
For the core certificate, the proposal-mass statement $\mu_K(G_\rho) \ge 1-2\rho$ is often sufficient in practice. When $C_\rho$ is small (say $< 1$), the importance weights on $G_\rho$ are nearly uniform, so $\bar{\pi}_z(G_\rho) \approx \mu_K(G_\rho) \approx 1-2\rho$. The target-mass conversion in Proposition~\ref{prop:target-mass} is a useful theoretical addition but is not essential for the core certificate's practical value.
\end{remark}


%% ═══════════════════════════════════════════════════════
\section{Dual-Certificate Framework}
%% ═══════════════════════════════════════════════════════

The V1 and V2 certificates serve complementary roles:

\medskip
\begin{center}
\begin{tabular}{lcc}
\hline
& \textbf{V1: Covering} & \textbf{V2: Quantile Core} \\
\hline
Region & Full $K_\alpha$ & Residual core $G_\rho$ \\
Controls & Worst-case oscillation & Core $(1-2\rho)$-mass oscillation \\
Rate & $O(n^{-1/D})$ & $O(n^{-1/2})$ \\
Rigorous for & $D \lesssim 5$ & All $D$ (on the core) \\
Spectral gap type & Full $K_\alpha$-restricted & Core $G_\rho$-restricted \\
\hline
\end{tabular}
\end{center}

\medskip
\noindent Together, they provide a complete picture:
\begin{itemize}[nosep]
  \item V1 certifies global correctness where tractable (low $D$).
  \item V2 certifies high-mass core mixing in all dimensions, honestly reporting the uncertified tail fraction $2\rho$.
  \item When both are available, V1 provides the stronger guarantee; V2 extends certification to regimes where V1 is vacuous.
\end{itemize}


%% ═══════════════════════════════════════════════════════
\section{Empirical Calibration (Preview)}
%% ═══════════════════════════════════════════════════════

Preliminary experiments on the $D$-dimensional banana target with CerT-OG-Anneal training give (with $\rho = 0.01$, $n = 200{,}000$):

\medskip
\begin{center}
\begin{tabular}{rrrrrr}
\hline
$D$ & $\osc_{K_\alpha}(r)$ & $\hat{C}_{0.01}$ & $\gamma_{\mathrm{V1}}$ & $\gamma_{\mathrm{core},0.01}$ & Ratio \\
\hline
2  & 2.27 & 0.13 & 0.152 & 0.933 & 6$\times$ \\
5  & 2.91 & 0.22 & 0.020 & 0.890 & 44$\times$ \\
6  & 2.50 & 0.19 & 0.010 & 0.908 & 94$\times$ \\
8  & 2.11 & 0.20 & 0.011 & 0.902 & 84$\times$ \\
10 & 2.27 & 0.27 & 0.003 & 0.865 & 292$\times$ \\
20 & 54.65 & 1.79 & $\approx 0$ & 0.286 & $\infty$ \\
\hline
\end{tabular}
\end{center}

\medskip
\noindent The core oscillation $\hat{C}_{0.01}$ remains between 0.13 and 0.27 for $D = 2$--$10$ (near dimension-free), while V1's covering certificate degrades from 0.15 to 0.003. At $D=20$, V1 is completely vacuous ($C_{\mathrm{V1}} = 96.2$), while the core certificate remains non-vacuous ($\hat{C}_{0.01} = 1.79$, $\gamma_{\mathrm{core}} = 0.29$). The increase at $D=20$ reflects genuine flow quality degradation, not a proof artifact.


%% ═══════════════════════════════════════════════════════
\section{Open Directions}
%% ═══════════════════════════════════════════════════════

\begin{enumerate}[label=(\alph*)]
  \item \textbf{From core to full spectral gap.} Theorem~\ref{thm:core-gap} certifies the restricted kernel $P_{G_\rho}$. Extending to the full $K_\alpha$-restricted kernel requires a comparison theorem relating $\gamma(P_{K_\alpha})$ and $\gamma(P_{G_\rho})$, accounting for the $2\rho$ tail mass. Possible routes: conductance comparison, lazy-chain domination, or $s$-conductance bounds.

  \item \textbf{Sub-Gaussian residuals.} If $r(Z)$ under $\mu_K$ can be shown to be sub-Gaussian with parameter $\sigma$, then $C_\rho \le 2\sigma\sqrt{2\log(1/\rho)}$, giving a parametric characterisation of core quality. The empirical data (variance $\approx 0.001$--$0.002$ for $D=2$--$10$) suggests sub-Gaussianity is plausible.

  \item \textbf{Certificate-aware training for V2.} The CerT-OG-Anneal loss minimises the full oscillation $\osc(r)$. For V2, a more natural training objective would directly minimise $\mathrm{Var}_{\mu_K}(r)$ or the $(1-2\rho)$-quantile range, which targets the core certificate rather than the worst-case tails.

  \item \textbf{Rigorous target-mass conversion.} Proposition~\ref{prop:target-mass} gives a bound involving either $\sup r$ or finite-sample exponential moment estimates. A self-normalised importance sampling concentration inequality (e.g., Chatterjee--Diaconis 2018) could give a clean finite-sample bound on $\bar{\pi}_z(G_\rho)$.
\end{enumerate}


%% ═══════════════════════════════════════════════════════
\section*{References}
%% ═══════════════════════════════════════════════════════

\begin{itemize}[nosep,leftmargin=*]
  \item Mengersen, K.L.\ and Tweedie, R.L.\ (1996). Rates of convergence of the Hastings and Metropolis algorithms. \emph{Annals of Statistics}, 24(1):101--121.
  \item Dvoretzky, A., Kiefer, J., and Wolfowitz, J.\ (1956). Asymptotic minimax character of the sample distribution function and of the classical multinomial estimator. \emph{Annals of Mathematical Statistics}, 27(3):642--669.
  \item Massart, P.\ (1990). The tight constant in the Dvoretzky--Kiefer--Wolfowitz inequality. \emph{Annals of Probability}, 18(3):1269--1283.
  \item Tierney, L.\ (1994). Markov chains for exploring posterior distributions. \emph{Annals of Statistics}, 22(4):1701--1728.
\end{itemize}

\end{document}
